\section*{Sources \& Limitations}
\addcontentsline{toc}{section}{Sources \& Limitations}

\subsection*{Primary sources consulted}
10月7日のOpenAIのGPT-6とIntelligent UI、10月7日のAnthropicのHaiku 5.5、10月6日のOpenAIのDecisions API betaと10月6日から7日のPerplexityの決定・埋め込み、10月8日のUltrafast tierとGemini仕事エージェント、10月3日のAleph AlphaのKolibriと10月7日のLiquidのd1、10月6日のEmbeddingGemma 2とReflectionのBeam（10月8日更新）とMistral Large 4、10月9日のCloudflareのClef-omniと10月5日のNova Sonic、10月6日のNano Banana 2.1カード、10月5日のOpenAIのtextGrainと10月7日のSynthID拡張、10月8日のAnthropicのCyber Missionと10月6日のOpenAIの数学公開・Ironclad模擬を範囲として読んだ。本号では全文PDFや手法の細目、図表の詳細までは扱わず、その範囲を超える主張は行わない。

\subsection*{Temporal boundary}
通常の範囲は2026-10-02 18:00から2026-10-09 18:00（いずれもAmerica/New\_York、終了時刻は含まない）。10月9日22時04分（UTC）の意図せざる動作報告と同日23時11分（UTC）のClef-omni拡張告知は範囲の区切りとともに読む。リポジトリの作成をもって重みの公開完了と読まない。日付は日の単位のものが多く、時刻レベルの断定は確かめたものに限る。

\subsection*{Vendor-reported boundary}
測定値や比較、コスト、スループット、利用見通しの説明は、ベンダーやモデルの配布元、提携の報告として受け止め、独立した再現とは切り分ける。GPT-6の44\%やHaikuの料金表とベンチ表、Decisionsの10倍表記、Perplexityの指数と単価、Ultrafastの提供範囲、KolibriやBeam、Mistralの規模とベンチ、Nano Bananaの横並び値、textGrainの検出率と品質表、SynthIDの規模、Cyberの見込み値、数学の計算統計、Ironcladの模擬相対値といった値はベンダー側の報告であり、条件の違いを含めて本号の検証範囲外である。

\subsection*{Community observation boundary}
公開投稿の台帳は原則として関心の方向を示す背景にとどめる。ただし、一次発表として本文で明示的に引用した公式アカウントの投稿は、その発表内容に限って一次資料として扱う。台帳に載っているという事実だけでは、いかなる行も技術的な根拠に格上げしない。独立した投稿やコミュニティの反応は、仕様やベンチマーク、日付、利用条件、価格、ライセンス、アーキテクチャ、公開事実の根拠にはならない。全行の内訳は本号の公開台帳で直接確かめられる。行単位の台帳を正とし、要約文の数えに頼らない。

\subsection*{What was not elevated}
持ち越しの4件は、新しい日付のある証拠がなく今号の章には載せない。いずれも対象期間内の素材的に重い品であることは認め、載せない理由を隠さない。DrexやStepFun、Decision-1等の高関心の持ち越しは未解決のまま残し、次号の取得対象とする。textGrain技術報告PDF、BeamとMistralの重み、system-card全文、料金・ライセンス条項も取得対象である。収集の台帳そのものは方法論の覚えに残し、物語には入れない。噂の段階の話題は今週の技術的な出来事として取り上げない。
